% problem-000292 1 钛铁矿提取金属钛
\section*{1. 钛铁矿提取金属钛}
（图：）

\subsection*{1-1}
钛铁矿的主要成分为 $\mathrm { F e T i O _ { 3 \mathrm { c } } }$ 。控制电炉熔炼温度(<1500K)，⽤碳还原出铁，⽽钛则进⼊炉渣浮于熔融铁之上，使钛与铁分离，钛被富集。已知：

$
\mathrm{FeTiO} _ {3} + \mathrm{C} \rightarrow \mathrm{TiO} _ {2} + \mathrm{Fe} + \mathrm{CO}
$

$
\Delta_ {\mathrm{r}} \mathrm{G} _ {\mathrm{m}} / \mathrm{J} \mathrm{mol} ^ {- 1} = 1 9 0 9 0 0 - 1 6 1 T
$

$
\mathrm{FeTiO} _ {3} + 4 \mathrm{C} \rightarrow \mathrm{TiC} + \mathrm{Fe} + 3 \mathrm{CO}
$

$
\Delta_ {\mathrm{r}} \mathrm{G} _ {\mathrm{m}} / \mathrm{J} \mathrm{mol} ^ {- 1} = 7 5 0 0 0 0 - 5 0 0 T
$

$
\mathrm{FeTiO} _ {3} + 3 \mathrm{C} \rightarrow \mathrm{Ti} + \mathrm{Fe} + 3 \mathrm{CO}
$

$
\Delta_ {\mathrm{r}} \mathrm{G} _ {\mathrm{m}} / \mathrm{J} \mathrm{mol} ^ {- 1} = 9 1 3 8 0 0 - 5 1 9 T
$

通过计算判断在电炉熔炼中主要发⽣以上哪个反应？

\subsection*{1-2}
写出在1073-1273 K下氯化反应的化学⽅程式。

\subsection*{1-3}
氯化得到的 $\mathrm { T i C l _ { 4 } }$ 中含有的 $\mathrm { V O C l _ { 3 } }$ 必须⽤⾼效精馏的⽅法除去，为什么？实际⽣产中常在409 K下⽤Cu还原 $\mathrm { { V O C l } _ { 3 } } ,$ ，反应物的摩尔⽐为1:1，⽣成氯化亚铜和难溶于 $\mathrm { T i C l _ { 4 } }$ 的还原物，写出还原反应⽅程式。

\subsection*{1-4}
精制后的 $\mathrm { T i C l _ { 4 } }$ ⽤⾦属镁还原可得海绵钛，写出化学反应⽅程式。

\subsection*{1-5}
菱镁矿（主要成分为 $\mathrm { M g C O _ { 3 } ) }$ 煅烧分解后与焦炭混合在氯化器中加热到1373 K，通⼊氯⽓⽣成$\mathrm { M g C l _ { 2 \circ } }$ 在 1023 K 电解熔融 $\mathrm { M g C l _ { 2 } }$ 得到⾦属镁。

\subsection*{1-5}
-1 写出⽤菱镁矿煅烧及氯化制取 $\mathrm { M g C l _ { 2 } }$ 的化学反应⽅程式。

\subsection*{1-5}
-2 写出电解熔融 $\mathrm { M g C l _ { 2 } }$ 的电极反应式和电解反应式。

\subsection*{1-5}
-3 已知 1023 K MgCl_{2}(l)的标准摩尔⽣成焓 $\Delta _ { \mathrm { f } } H _ { \mathrm { m } } ^ { \ominus }$ 为−596.32 kJ $\mathrm { m o l ^ { - 1 } }$ ，MgCl_{2}(l)、 $\mathrm { { M g ( l ) } }$ $\mathrm { C l _ { 2 } ( g ) }$ 的标准摩尔熵 $S _ { \mathrm { m } } ^ { \ominus }$ 分别为 231.02、77.30、268.20 kJ $\mathrm { m o l ^ { - 1 } } _ { \mathrm { c } }$ 。计算MgCl_{2}的理论分解电压。
